% INSERCIÓN FOCAL AUTORIZADA: no introduce un capítulo ni el concepto de radión.
% Incluir dentro de la sección constitutiva, después de la inversión cúbica
% y de la declaración de las cartas dimensionales. Etiquetas fuera de grupos
% de prefijado automático. Conserva íntegramente el bloque geométrico anterior.
\subsection{Reconstruction of the action shape from the constitutive response}
\label{iii:subsec:puente-forma-LC}

The constitutive response and the inductive--capacitive realization receive
the same angular pair constructed from APP--TRIT--TPK. The former evaluates
rational functions of the two channels; the latter uses the ratio of
their exponents to determine the shape of the action ellipse. Cubic
inversion allows the two readouts to be composed without identifying their
numerical ratios.

Let $x=A_{\rm rad}$ and $y=C^*_{\rm rad}$, with $x>|y|$, be the lifted
coordinates already constructed, and preserve their sheet labels. Write
\begin{equation}
 q_\pm=e^{-(x\pm y)},\qquad s_\pm=q_\pm^{30},\qquad
 r_\pm=f(s_\pm),\qquad
 f(s)=\frac{1+s+s^2}{1+s+s^2+s^3}.
 \label{iii:eq:forma-canales}
\end{equation}
The contractive domain gives $s_\pm\in(0,1)$ and
$r_\pm\in(3/4,1)$. In the chamber $x>y>0$, we have
$s_+<s_-$ and $r_+>r_-$; sheet interchange transports this chamber
to the opposite orientation.

The image of the two normalized constitutive coordinates is the domain
\begin{equation}
 \mathcal U=\left\{(z,v)\in(0,\infty)^2:
   (zv)^{-1/2}\in(3/4,1),\quad
   (z/v)^{1/2}\in(3/4,1)\right\}.
 \label{iii:eq:forma-dominio}
\end{equation}
Here $(z,v)=(\widehat Z,\widehat c)$. The restrictions express
membership of the reconstructed channels in the domain of the inversion
theorem; they do not constitute a subsequent choice of their values.

\Needspace{23\baselineskip}
\begin{proposition}[Constitutive composition of the shape and LC impedance]
\label{iii:prop:forma-LC}
For $(\widehat Z,\widehat c)\in\mathcal U$, let $s_\pm\in(0,1)$ be the
unique roots of
\begin{equation}
 r_\pm s_\pm^3+(r_\pm-1)(1+s_\pm+s_\pm^2)=0,
 \qquad
 r_+=\frac1{\sqrt{\widehat Z\widehat c}},\quad
 r_-=\sqrt{\frac{\widehat Z}{\widehat c}}.
 \label{iii:eq:forma-inversion}
\end{equation}
The angular ratio, the anisotropy of the action ellipse, and the relative LC
impedance are reconstructed exactly by
\begin{align}
 \frac{y}{x}
 &=\frac{\log s_+-\log s_-}{\log s_++\log s_-},
 \label{iii:eq:forma-razon}\\
 \eta_{\rm el}
 &=\frac12\log\frac{\log s_+}{\log s_-},
 \label{iii:eq:forma-eta}\\
 \frac{Z_{\rm LC}}{Z_*}
 &=e^{-\eta_{\rm el}}
  =\sqrt{\frac{\log s_-}{\log s_+}}.
 \label{iii:eq:forma-LC}
\end{align}
In the last identity, $Z_{\rm LC}=Z_{\eta_{\rm el}}$ is the impedance
of the LC chart in \eqref{eq:ii-elipse-lc}.
\end{proposition}
\begin{proof}
The cubic inversion in \eqref{iii:eq:cubic} uniquely determines $s_\pm$,
because $f$ is a decreasing bijection from $(0,1)$ onto
$(3/4,1)$. Their positive roots of order $30$ recover $q_\pm$.
The identities
\[
 \log s_+=-30(x+y),\qquad \log s_-=-30(x-y)
\]
prove \eqref{iii:eq:forma-razon} by difference and sum. Both logarithms
are negative; their ratio is positive and their sum is nonzero. Therefore,
\[
 \frac12\log\frac{\log s_+}{\log s_-}
   =\frac12\log\frac{x+y}{x-y}
   =\operatorname{artanh}(y/x)=\eta_{\rm el}.
\]
The relation $Z_{\eta_{\rm el}}/Z_*=e^{-\eta_{\rm el}}$, proved in
\eqref{eq:ii-elipse-impedancia}, gives \eqref{iii:eq:forma-LC}.
\end{proof}

The composition preserves the information needed for each reconstruction.
With $\hbar$, one recovers
$a_S=\sqrt{2\hbar}\,e^{\eta_{\rm el}/2}$ and
$b_S=\sqrt{2\hbar}\,e^{-\eta_{\rm el}/2}$; with the chart
$(\omega,Z_*)$, one recovers the inductance and capacitance already defined.
Action, clock, and dimensional reference accompany the shape as
typed data. Inverting two dimensionless coordinates does not
reconstruct their units or the entire memory of the state.

Normalization enters before the cubic equations are solved.
If $Z_{\rm vac}$ and $c_{\rm vac}$ are received in a dimensional chart,
relations \eqref{iii:eq:dimensions} first determine
\[
 \widehat Z=\frac{Z_{\rm vac}}{u_Su_Q^{-2}},\qquad
 \widehat c=\frac{c_{\rm vac}}{u_C}.
\]
These are the coordinates to which
\eqref{iii:eq:forma-inversion} applies. The reference $Z_*$ belongs to the
LC realization and is explicitly preserved in \eqref{iii:eq:forma-LC}.

The two ratios are related through the channels, but evaluate
different functions:
\begin{equation}
 \widehat Z=\frac{f(s_-)}{f(s_+)},\qquad
 \frac{Z_{\rm LC}}{Z_*}
       =\sqrt{\frac{\log s_-}{\log s_+}}.
 \label{iii:eq:forma-dos-cocientes}
\end{equation}
The preceding reconstruction provides the map between the two readouts.
In particular, interchanging the sheets inverts both ratios, changes
$\eta_{\rm el}$ to $-\eta_{\rm el}$, and preserves $\widehat c$.
In the symmetric limit $y=0$, the channels coincide and both ratios
equal $1$. Outside that case, a common origin does not impose equality
of the two functions. The action shape is recoverable from the
constitutive response while preserving orientation and normalization.
